\section{DSPy: programación declarativa de IA}

\subsection{Idea central}

\begin{importantbox}{La idea central de DSPy}
DSPy permite que los desarrolladores definan la estructura y la lógica de un sistema de IA mediante un \textbf{programa de Python}, y el marco de trabajo \textbf{compila} automáticamente el prompt y los ejemplos few-shot óptimos. Los desarrolladores solo necesitan declarar «qué hacer» (what); no necesitan especificar manualmente «cómo hacerlo» (how, es decir, el texto concreto del prompt).
\end{importantbox}

Abstracciones clave:
\begin{itemize}
    \item \textbf{Signature}: define los tipos de entrada y salida de un módulo (por ejemplo, «question -> answer»)
    \item \textbf{Module}: componente de IA componible (como ChainOfThought, ReAct)
    \item \textbf{Optimizer}: busca automáticamente el prompt, los ejemplos few-shot y la configuración de módulos óptimos
    \item \textbf{Metric}: define el criterio de evaluación de la calidad del sistema
\end{itemize}

\subsection{Optimizadores (Optimizers)}

DSPy ofrece varios optimizadores para mejorar el sistema de manera automática:

\begin{itemize}
    \item \textbf{BootstrapFewShot}: selecciona automáticamente los mejores ejemplos few-shot a partir de los datos de entrenamiento
    \item \textbf{COPRO}: optimiza automáticamente el texto de instrucción dentro del prompt
    \item \textbf{MIPRO}: optimiza simultáneamente las instrucciones y los ejemplos
    \item \textbf{BootstrapFinetune}: genera datos de entrenamiento y aplica fine-tuning al modelo subyacente
\end{itemize}

\subsection{Compilación frente a escritura manual}

\begin{knowledgebox}{Analogía con la programación tradicional}
DSPy es a prompt engineering lo que un compilador es al lenguaje ensamblador. El desarrollador escribe el programa con abstracciones de alto nivel, y el compilador (optimizer) genera automáticamente la implementación subyacente (el prompt). Esto no solo mejora la eficiencia, sino que también dota al sistema de portabilidad: para cambiar de modelo basta con volver a compilar.
\end{knowledgebox}

\subsection{Resumen de la sección}

Mediante la programación declarativa y la optimización automática, DSPy eleva la construcción de sistemas de IA compuestos, de la prompt engineering manual, a una ingeniería de software sistemática.

